%----------头部信息-----------------
\begin{tabular*}{\textwidth}{l@{\extracolsep{\fill}}r}
  \textbf{\Large 夏驭风（Yufeng Xia）} & 手机：\href{tel:+8613862417657}{(+86) 138-6241-7657} \\
  \href{https://xyf2002.github.io/}{xyf2002.github.io} \quad
  \href{https://www.linkedin.com/in/yufeng-xia-12648a25b}{linkedin.com/in/yufeng-xia} & 微信：xyf2316639449 \\
  英国爱丁堡 & 邮箱：\href{mailto:yufeng.xia@ed.ac.uk}{yufeng.xia@ed.ac.uk} \\
\end{tabular*}

\vspace{2pt}

%-----------求职意向与简介-----------------
\section{求职意向与个人简介}
  \begin{itemize}[leftmargin=*]
    \skillrow{求职意向}{AI infra 与芯片架构工程师（软件研发类）· 工作地上海 · 2026 年 11 月获硕士学位（2027 届）· 2026 年 9 月起可实习或全职入职。}
    \item\small{爱丁堡大学计算系统架构研究所（ICSA）研究型硕士，导师 Mahesh K. Marina 教授，研究方向为 \textbf{AI 算力基础设施、加速器调度与性能工程}：GPU 算力共享与抢占、算力划分与干扰刻画、分布式训练系统的性能建模与大规模评测。
    \par\vspace{2pt}
    以第一或合作作者在 \textbf{MobiCom、EdgeSys（最佳论文）、SEC} 等会议发表系统方向工作；工程侧长期用 \textbf{C/C++} 编写事件驱动的分布式后端与容错逻辑，具备 \textbf{Triton GPU 算子开发}、\textbf{CUDA Green Context 算力分区}、\textbf{Verilog/FPGA 处理器实现}、\textbf{LLVM 编译优化} 与 \textbf{Linux 内核构建} 的完整链路经验，覆盖从 RTL 到框架的各层抽象。\vspace{-2pt}}
  \end{itemize}
